\documentclass[11pt]{article}
\usepackage[margin=1.1in]{geometry}
\usepackage{amsmath, amssymb, amsthm, mathtools}
\usepackage{booktabs}
\usepackage{enumitem}
\usepackage{hyperref}
\usepackage{parskip}
\usepackage{xcolor}
\hypersetup{colorlinks=true, urlcolor=blue, linkcolor=black}

\title{Prior accumulator models in unified notation\\
\large{Translating McMurray, Mitchell--McMurray, Hidaka, Mollica--Piantadosi,\\
       and Kachergis into our $(\alpha_i, r_i, \kappa_i, \delta_j)$ coordinates}}
\author{Standard Model 2 working note}
\date{\today}

\newcommand{\E}{\mathbb{E}}
\newcommand{\Var}{\mathrm{Var}}

\begin{document}
\maketitle

\section{Why this exercise}

Each prior accumulator model uses its own notation: McMurray writes a
threshold density $g(\tau)$, Hidaka uses rate-change exponent $D$,
Mollica \& Piantadosi use a per-word $(k_w, \lambda_w)$, and so on.
Translating each model into our common parameterization makes the
comparison literal rather than structural. It also pins down exactly
which of our parameters each prior model pins (and at what value),
versus what they leave free.

The unified parameterization throughout this note is
\begin{equation}
\eta_{ij}(t) \;=\;
\underbrace{\log r_i}_{\text{input rate}}
\;+\;
\underbrace{\log\alpha_i}_{\text{efficiency}}
\;+\;
\underbrace{\kappa_i\,\log t}_{\text{scaling}}
\;-\;
\underbrace{\delta_j}_{\text{word difficulty}},
\qquad
P(\text{produces}_{ij}\mid t) = \sigma(\eta_{ij}).
\label{eq:our-model}
\end{equation}
A pure accumulator pins $\alpha_i \equiv 1$ and $\kappa_i \equiv 1$,
leaving only $r_i$ and $\delta_j$ free. Our model frees both
$\alpha_i$ and $\kappa_i$ at the per-child level.

Each section below: the original model, in its native notation; the
translation; what's pinned vs.\ free in our coordinates.

\section{McMurray (2007): pure accumulator}

\paragraph{Original.}
Each word $w$ has a fixed acquisition threshold $\tau_w$, drawn from
a population distribution $g(\tau)$. At each unit time step every
word's accumulator gains one count. Word $w$ is learned at the first
$t$ such that the count reaches $\tau_w$. Since accumulation is
unit-rate and deterministic, that's the first $t \ge \tau_w$.
Aggregate vocabulary is $L(t) = \int_0^t g(\tau)\,d\tau = G(t)$
(the CDF of $g$).

\paragraph{Translation.}
A deterministic step $\mathbb{1}[t \ge \tau_w]$ is the zero-noise
limit of $\sigma\!\big(\lambda(t - \tau_w)\big)$ as $\lambda \to \infty$.
With our IRT logit link,
\begin{equation}
P(\text{produces}_w \mid t)
\;=\;
\sigma\!\big(\log t - \log \tau_w\big)
\;=\;
\sigma\!\big(\kappa\,\log t - \delta_w\big)
\quad\text{with}\quad
\kappa \equiv 1,\;\; \delta_w := \log\tau_w.
\end{equation}
The threshold density $g(\tau)$ becomes the prior over the per-word
$\delta_j$ on a log-time scale.

\paragraph{In our parameterization.}

\begin{center}
\renewcommand{\arraystretch}{1.10}
\begin{tabular}{@{}lll@{}}
\toprule
\textbf{Parameter} & \textbf{Setting in McMurray} & \textbf{Free in our model?} \\
\midrule
$\log r_i$         & $\equiv 0$ (constant rate, no per-child variation) & free \\
$\log \alpha_i$    & $\equiv 0$ (no efficiency variation) & free \\
$\kappa_i$         & $\equiv 1$ (unit-rate accumulator) & free \\
$\delta_j$         & $= \log \tau_j$, drawn from $g$ & free \\
\bottomrule
\end{tabular}
\end{center}

\paragraph{Acceleration in this model.} The aggregate vocabulary curve
$L(t) = G(t)$ is the CDF of $g$. Whether the population-mean trajectory
shows visible acceleration depends entirely on the shape of $g$:
\begin{itemize}\setlength\itemsep{0.2em}
\item If $g(\tau)$ is uniform on a window, $L(t)$ is linear in $t$
  over that window.
\item If $g'(\tau) > 0$ over the relevant range (more words have
  higher thresholds than lower), $L(t)$ accelerates.
\item If $g'(\tau) < 0$, $L(t)$ decelerates.
\end{itemize}
McMurray's argument: empirical $L(t)$ accelerates, so $g'(\tau) > 0$.
\textbf{This is a population-level claim, not a per-child claim.}
McMurray has no $\kappa_i$ heterogeneity ($\sigma_\kappa = 0$); the
``acceleration'' he refers to is in $g$, which corresponds to
heterogeneity in our $\delta_j$ across words.

\section{Mitchell \& McMurray (2009): McMurray + leverage}

\paragraph{Original.}
McMurray (2007) plus a feedback term: as the child accumulates
vocabulary, each known word ``leverages'' the acquisition of new
words. The effective time becomes
\begin{equation}
a(t) = t + c\,L(t),
\quad
L(t) = G(a(t)),
\quad
\frac{dL}{dt} = \frac{g(t + cL)}{1 - c\,g(t+cL)},
\end{equation}
where $c$ is the leverage strength.

\paragraph{Translation.}
The leverage $cL(t)$ couples each child's effective time to their own
vocabulary state. In our linear predictor, this is a non-linear-in-$L$
modification:
\begin{equation}
\eta_{ij}(t) = \log a(t) - \delta_j = \log\!\big(t + c\,L_i(t)\big) - \delta_j.
\end{equation}
The argument is no longer linear in $\log t$, and $a(t)$ depends on
the child's own vocabulary state, which makes per-(child, word) likelihoods
non-separable in our framework. \textbf{This is the one prior model
that doesn't embed cleanly in
equation~\eqref{eq:our-model}.}

\paragraph{Mitchell \& McMurray's result.} They prove the leverage
\emph{shifts} the growth curve forward or backward in time but
\emph{cannot create acceleration not already present in $g(\tau)$}.
In our coordinates: leverage moves the timing of acquisition without
changing $\kappa_i$ — it can't manufacture super-linear scaling.

\section{Hidaka (2013): closed-form AoA families}

\paragraph{Original.}
Three nested learning processes, each with closed-form AoA distribution:

\begin{center}
\renewcommand{\arraystretch}{1.10}
\begin{tabular}{@{}lll@{}}
\toprule
\textbf{Process} & \textbf{Dynamics} & \textbf{AoA distribution} \\
\midrule
Cumulative & $N$-accumulator at constant rate $\delta$ & $\Gamma(N, \delta)$ \\
Rate-change & 1-accumulator at rate $\delta \cdot t^D$ & Weibull \\
Cumulative + rate-change & both & Weibull--Gamma \\
\bottomrule
\end{tabular}
\end{center}

The rate-change form yields ability $\propto t^{D+1}$, since
$\int_0^t \delta \tau^D\,d\tau = \frac{\delta}{D+1}\,t^{D+1}$.

\paragraph{Translation: rate-change form.}
On the logit scale,
\begin{equation}
\log(\text{ability at } t) = (D+1)\log t + \mathrm{const},
\end{equation}
so equation~\eqref{eq:our-model} reduces to
\begin{equation}
\eta_{ij}(t) = (D+1)\log t - \delta_j
\quad\text{with}\quad \kappa_i \equiv D+1.
\end{equation}
\textbf{Hidaka's $D$ is exactly $\kappa - 1$ in our notation.} His
$D > 0$ is our $\kappa_{\text{pop}} > 1$ (super-linear scaling).
Crucially, $D$ is shared across all kids — no per-child variation.

\paragraph{Translation: cumulative form.}
The Gamma waiting-time form $\Gamma(N_w, \delta)$ does not have a
simple closed-form on the log scale; for large $N_w$ it converges
to a normal CDF, which approximates a sigmoid. Mapping at the
mean: $\delta_j \approx \log(N_j/\delta)$, and $\kappa_i \equiv 1$.
This is essentially McMurray's pure accumulator with a smoothed
threshold.

\paragraph{In our parameterization (rate-change).}

\begin{center}
\renewcommand{\arraystretch}{1.10}
\begin{tabular}{@{}lll@{}}
\toprule
\textbf{Parameter} & \textbf{Setting in Hidaka rate-change} & \textbf{Free in our model?} \\
\midrule
$\log r_i$         & $\equiv \log \delta_{\text{const}}$ (population rate) & free \\
$\log \alpha_i$    & $\equiv 0$ (no efficiency variation) & free \\
$\kappa_i$         & $\equiv D + 1$ (shared, no per-child variation) & free \\
$\delta_j$         & per-word, free                                  & free \\
\bottomrule
\end{tabular}
\end{center}

Hidaka's posit: $D$ is informative about word properties.
He reports correlation between $D$ and imageability, taking $D$
from population-pooled fits. In our coordinates this corresponds
to per-word $\kappa_w$ (not per-child $\kappa_i$), which we don't
currently estimate. (It's an interesting direction.)

\section{Mollica \& Piantadosi (2017): per-word Gamma waiting time}

\paragraph{Original.}
Per-word generative model: word $w$ acquisition begins at start age
$s_w$, ``effective learning instances'' (ELIs) arrive as Poisson at
rate $\lambda_w$, and the word is acquired after $k_w$ ELIs. So
\begin{equation}
\mathrm{AoA}_w \sim s_w + \Gamma(k_w, \lambda_w),
\quad
F_w(a) = \int_0^{a-s_w} \frac{u^{k_w-1}\,e^{-u\lambda_w}\,\lambda_w^{k_w}}{\Gamma(k_w)}\,du.
\end{equation}
Mollica \& Piantadosi report $k \approx 10$, $\lambda \approx 0.5$/mo,
$s \approx 2$~mo for comprehension across 13 languages.

\paragraph{Translation.}
For large $k_w$, the Gamma CDF converges to a Gaussian CDF (CLT) with
mean $k_w/\lambda_w$ and standard deviation $\sqrt{k_w}/\lambda_w$.
The Gaussian CDF approximates the logistic with a $\sim$1.7 scale
factor. So
\begin{equation}
F_w(a) \approx \sigma\!\big(b_w (a - k_w/\lambda_w)\big),
\quad
b_w \propto \lambda_w/\sqrt{k_w}.
\end{equation}
On the log scale (writing $a = e^{\log t}$, no $s_w$ for simplicity),
the curve is approximately a sigmoid in $\log t$ with slope
proportional to $b_w \cdot k_w/\lambda_w = \sqrt{k_w}$.

In our parameterization:
\begin{equation}
\eta_{ij}(t) = \log t - \delta_j
\quad\text{with}\quad
\delta_j \approx \log(k_j/\lambda_j),\;\;
\kappa_i \equiv 1.
\end{equation}
$\kappa$ doesn't appear directly because the Gamma waiting time has
a fixed shape; the per-word slope variation comes from $\sqrt{k_w}$
inside the sigmoid scale, which our model doesn't have a free
parameter for (we'd need a per-word discrimination $\lambda_j$, the
2PL extension).

\paragraph{In our parameterization.}

\begin{center}
\renewcommand{\arraystretch}{1.10}
\begin{tabular}{@{}lll@{}}
\toprule
\textbf{Parameter} & \textbf{Setting in Mollica \& Piantadosi} & \textbf{Free in our model?} \\
\midrule
$\log r_i$         & $\equiv \log \lambda$ (shared rate) & free \\
$\log \alpha_i$    & $\equiv 0$ (no per-child variation) & free \\
$\kappa_i$         & $\equiv 1$ (unit-rate accumulator) & free \\
$\delta_j$         & $= \log(k_j/\lambda_j)$, free       & free \\
\bottomrule
\end{tabular}
\end{center}

\section{Kachergis et al.\ (2021): 1-PL Rasch IRT with covariates}

\paragraph{Original.}
\begin{equation}
\mathrm{logit}\,P(y_{ij}=1) = \theta_i + d_j + \beta_a\cdot a_i + \beta_t\cdot \log(p_j H a_i),
\end{equation}
where $\theta_i$ is per-child latent ability, $d_j$ is per-word
difficulty, $a_i$ is age, $p_j$ is the word's CHILDES probability,
$H$ is waking hours per month, and $\beta_a, \beta_t$ are shared
slopes (no per-child variation).

\paragraph{Translation.}
Identify $\theta_i = \log r_i + \log\alpha_i$ (their per-child
intercept), $d_j = -\delta_j$ (their per-word intercept), and the
linear age covariate $\beta_a a_i$ as a linearization of
$\kappa\,\log a$:
\begin{equation}
\eta_{ij} \;\approx\;
(\log r_i + \log\alpha_i) + \log p_j + \log H + \kappa\,\log a_i - \delta_j,
\end{equation}
with $\kappa \equiv \beta_a/(\partial \log a /\partial a) \cdot a$ a
constant population slope, no per-child variation.

\paragraph{In our parameterization.}

\begin{center}
\renewcommand{\arraystretch}{1.10}
\begin{tabular}{@{}lll@{}}
\toprule
\textbf{Parameter} & \textbf{Setting in Kachergis 2021} & \textbf{Free in our model?} \\
\midrule
$\log r_i$         & free (population $\mu_r$, no per-child variance specified) & free, $\sigma_r$ pinned \\
$\log \alpha_i$    & free (their $\theta_i$ absorbs $\log r + \log \alpha$) & free \\
$\kappa_i$         & $\equiv$ shared population slope (no per-child variance) & free \\
$\delta_j$         & free (their $d_j$)                  & free \\
\bottomrule
\end{tabular}
\end{center}

\textbf{Kachergis et al.\ are the first prior model to free $\sigma_\alpha$.}
But $\sigma_\kappa$ is still zero: every child has the same
acceleration slope.

\section{Unified table}

In our coordinates, every prior model pins or frees specific parameters.
The pattern below makes clear what's structural and what's notational
across the literature.

\begin{center}
\renewcommand{\arraystretch}{1.20}
\begin{tabular}{@{}lcccc@{}}
\toprule
\textbf{Model} &
$\boldsymbol{\sigma_\alpha}$ &
$\boldsymbol{\kappa_{\text{pop}}}$ &
$\boldsymbol{\sigma_\kappa}$ &
\textbf{Per-word $\delta_j$?} \\
\midrule
McMurray (2007)                & $0$        & $1$       & $0$ & free, $\delta_j = \log\tau_j$ \\
Mitchell \& McMurray (2009)    & $0$        & $1^{*}$   & $0$ & free, plus leverage $c$ \\
Hidaka (2013) cumulative       & $0$        & $1$       & $0$ & free, $\delta_j = \log(N_j/\delta)$ \\
Hidaka (2013) rate-change      & $0$        & $D{+}1$   & $0$ & free \\
Mollica \& Piantadosi (2017)   & $0$        & $1$       & $0$ & free, $\delta_j = \log(k_j/\lambda_j)$ \\
Kachergis et al.\ (2021)       & free       & free      & $0$ & free \\
\midrule
\textbf{Best-fitting model (this work)} &
\textbf{free} & \textbf{free} & \textbf{free} & \textbf{free} \\
\bottomrule
\end{tabular}
\end{center}

\noindent\footnotesize{$^{*}$Mitchell \& McMurray's leverage is not
linear-in-$\log t$, so $\kappa$ as defined in our model isn't quite
the right summary. Their model lives outside our power-law family.}

\section{What's structurally new in our model}

Reading down the table, the prior literature progressively frees
parameters: McMurray gives all four to a single shared form;
Hidaka's rate-change frees $\kappa_{\text{pop}}$; Kachergis frees
$\sigma_\alpha$. Each step adds one degree of freedom.

\textbf{Our model is the first to free $\sigma_\kappa$:} per-child
variation in the scaling exponent. That's the load-bearing addition
that enables fitting the variability signature jointly with
acceleration. The 4-panel comparison (deck slide 13) shows visually
that this addition is necessary --- variants without
$\sigma_\kappa$ free can't reproduce both the percentile fan width
and its growth with age.

\section{Discussion: what this implies for the literature}

\paragraph{Acceleration was always in the field, but as a
population-level claim.} McMurray's $g'(\tau) > 0$ argument and
Hidaka's $D > 0$ both posit acceleration at the population level.
We replicate that claim ($\kappa_{\text{pop}} \approx 10.4 \gg 1$
in English M\_best). What's new is per-child variance: kids differ
substantially in their scaling exponent.

\paragraph{Variability was treated as a measurement issue rather
than a structural feature.} Kachergis's $\theta_i$ first frees the
intercept, but the slope is shared. Treating between-child variance
as primarily intercept-level is consistent with much of the field's
focus on ``per-child ability'' as a single number. The data argue
for a richer two-axis structure (intercept + slope variance).

\paragraph{The per-child slope variance is a falsifiable claim, not
a parameter convenience.} If $\sigma_\kappa$ were small, our model
would collapse onto Kachergis (2021). The empirical posterior gives
$\sigma_\kappa = 3.47$ \,[3.14, 3.84] in English, well above zero
and large compared to $\sigma_\alpha = 1.71$. This is the
quantitative finding that requires explanation.

\end{document}
